\documentclass[12pt,letterpaper]{article}



%%
%% Required packages:
%%
\usepackage{etex}
\usepackage{amsmath}
\usepackage{amssymb}
\usepackage{amstext}
\usepackage{amsthm}
\usepackage{fancyhdr,fancybox}
\usepackage{mathrsfs} % need for \mathscr command
\usepackage{multicol}
\usepackage[normalem]{ulem}
%\usepackage{pictex,pictexplus}
% \usepackage{pictexwd}
\usepackage{rotating}
\usepackage{url}
\usepackage{xfrac}
\usepackage{booktabs}
\usepackage{color,soul}
\usepackage{comment}

\usepackage{hyperref}
\hypersetup{
    colorlinks=true,     % false: boxed links; true: colored links
    linkcolor=blue,      % color of internal links
    citecolor=blue,      % color of links to bibliography
    filecolor=blue,      % color of file links
    urlcolor=blue        % color of external links
}


\usepackage{tikz}
\usetikzlibrary{backgrounds,calc}

%%
%% Common formatting commands
%%
\setlength{\parindent}{0in}
\setlength{\textwidth}{7in}
\setlength{\evensidemargin}{-0.25in}
\setlength{\oddsidemargin}{-0.25in}
\setlength{\parskip}{.5\baselineskip}
\setlength{\topmargin}{-0.5in}
\setlength{\textheight}{9in}

%               Problem and Part
\newcounter{problemnumber}
\newcounter{partnumber}
\newcounter{subpartnumber}
\newcommand{\Problem}{%
  \refstepcounter{problemnumber}%
  \setcounter{partnumber}{0}%
  \setcounter{subpartnumber}{0}%
  \item[\makebox{\hfil\textbf{\theproblemnumber.}\hfil}]%
  }
\newcommand{\InProblem}[2][1.6]{%
  \refstepcounter{problemnumber}%
  \setcounter{partnumber}{0}%
  \setcounter{subpartnumber}{0}%
  \textbf{\theproblemnumber.}\ \ \parbox[t]{#1in}{#2}%
}
\newcommand{\InSmallProblem}[1]{\InProblem[1.05]{#1}}
\newcommand{\NoProblem}[1]{%
  \mbox{\hfil\textbf{#1}\hfil}%
  }
\newcommand{\UnProblem}{%
  \refstepcounter{problemnumber}%
  \setcounter{partnumber}{0}%
  \setcounter{subpartnumber}{0}%
  \mbox{\hfil\textbf{\theproblemnumber.}\hfil}%
  }
\newcommand{\InFlexProblem}[1]{%
  \refstepcounter{problemnumber}%
  \setcounter{partnumber}{0}%
  \setcounter{subpartnumber}{0}%
  \textbf{\theproblemnumber.}\ \ #1%
}
%%  Part:
\newcommand{\Part}{%
  \renewcommand{\thepartnumber}{(\alph{partnumber})}%
  \refstepcounter{partnumber}
  \setcounter{subpartnumber}{0}%
  \item[(\alph{partnumber})]%
}
\newcommand{\InPart}[2][1.75]{%
  \renewcommand{\thepartnumber}{(\alph{partnumber})}%
  \refstepcounter{partnumber}%
  \setcounter{subpartnumber}{0}%
  (\alph{partnumber})\ \ \parbox[t]{#1in}{#2}%
}
\newcommand{\UnPart}{%
  \refstepcounter{partnumber}%
  \renewcommand{\thepartnumber}{(\alph{partnumber})}%
  \setcounter{subpartnumber}{0}%
  (\alph{partnumber})%
}
\newcommand{\InLargePart}[1]{\InPart[3.05]{#1}}
\newcommand{\InFullPart}[1]{\InPart[6.2]{#1}}
\newcommand{\InSmallPart}[1]{\InPart[1.05]{#1}}
%% SubPart:
\newcommand{\SubPart}{%
  \refstepcounter{subpartnumber}%
  \item[(\roman{subpartnumber})]%
  }
\newcommand{\InSubPart}[2][2.25]{%
  \stepcounter{subpartnumber}%
  (\roman{subpartnumber})\ \ \parbox[t]{#1in}{#2}%
}
\newcommand{\InSmallSubPart}[1]{\InSubPart[1.5]{#1}}
\newcommand{\InTinySubPart}[1]{%
  \stepcounter{subpartnumber}%
  (\roman{subpartnumber})\ \ \makebox{#1}%
}
\newcommand{\InMySubPart}[2][2.25]{%
  \stepcounter{subpartnumber}%
  \parbox[t]{#1in}{\makebox[0.3in][l]{(\roman{subpartnumber})}\ \ {#2}}%
}
\newcommand{\TrueFalse}{\textbf{T}\ \ \textbf{F}\ \ }
\newcommand{\TFPart}[1]{% 
  \stepcounter{partnumber}%
  \setcounter{subpartnumber}{0}%
  \item[(\alph{partnumber})] \TrueFalse\parbox[t]{5.5in}{#1}}

\newcommand{\Points}[1]{(#1 points) \ }
\newcommand{\Point}[1]{(#1 point) \ }


%% For Answers:
\newcounter{answernumber}
\newcommand{\Answer}{%
  \refstepcounter{answernumber}%
  \setcounter{partnumber}{0}%
  \setcounter{subpartnumber}{0}%
  \item[\makebox{\hfil\textbf{\theanswernumber.}\hfil}]%
  }


% Setting up the headers for the first page
%\chead{} % will default to what is typed in the file.
\fancyhead[L]{\textsc{Math 22a}}
\fancyhead[R]{\textsc{Fall 2024}}
\fancyfoot[R]{
	\vspace*{\fill}\footnotesize\textcopyright{} \emph{Harvard Math 22a}	
}
\renewcommand{\headrulewidth}{0pt}

%\fancyhead[C]{}
%	\chead{}	
%	\lhead{}
%	\rhead{}
%	\cfoot{}


% Putting copyright on the footer of each page
%\fancypagestyle{secondpage}{
%	\fancyhead[C]{TESTing again} % change this to be blank
%		%\chead{}	
%	\fancyhead[L]{bears} % change this to be blank
%	\fancyhead[R]{dogs} % change this to be blank
%	\fancyfoot[R]{ %keeps the footer the same
%		\vspace*{\fill}\footnotesize\textcopyright{} \emph{Harvard Math 22a}	
%	}
%}
%\renewcommand{\headrulewidth}{0pt}
%\pagestyle{fancy}
\pagestyle{empty}


%\lhead{}
%\rhead{}
%\rfoot{\vspace*{\fill}\footnotesize\textcopyright{} \emph{Harvard Math 22a}}
%\renewcommand{\headrulewidth}{0pt}
%\pagestyle{fancy}




%%
%% Common shortcut commands:
%%
\newcommand{\ds}{\displaystyle}
\newcommand{\sgn}{\operatorname{sgn}}
\renewcommand{\Pr}{\operatorname{P}}
\newcommand{\CPr}[3][{}]{\Pr_{\text{#1}}\left(#2\,|\,#3\right)}

\newcommand{\C}{\mathbb{C}}
\newcommand{\R}{\mathbb{R}}
\newcommand{\Z}{\mathbb{Z}}
\newcommand{\N}{\mathbb{N}}
\newcommand{\Q}{\mathbb{Q}}
\newcommand{\D}{\mathbb{D}}

\newcommand{\rref}{\operatorname{rref}}
\newcommand{\im}{\operatorname{im}}
\newcommand{\rank}{\operatorname{rank}}
\newcommand{\diag}{\operatorname{diag}}
\newcommand{\Span}{\operatorname{span}}
%\renewcommand{\vec}[1]{\mathbf{#1}}
\newcommand{\charpoly}{\mathscr{P}}
\newcommand{\nicefrac}[2]{\sfrac{#1}{#2}} % using xfrac package
\newcommand{\topology}{\mathcal{T}}
\newcommand{\Inn}{\operatorname{Inn}}

\newcommand{\blankbox}[2]{\fbox{\rule{#1}{0in}\rule{0in}{#2}}}

\newenvironment{myindentpar}[1]%
 {\begin{list}{}%
         {\setlength{\leftmargin}{#1}
          \setlength{\rightmargin}{\leftmargin}}%
         \item[]%
 }
 {\end{list}}
\newtheorem{theorem}{Theorem}
\newtheorem*{NStheorem}{Theorem}
\newtheorem{cor}[theorem]{Corollary}
\newtheorem*{claim}{Claim}
\newtheorem{defn}{Definition}

\newenvironment{amatrix}[1]{%
  \left(\begin{array}{@{}*{#1}{c}|c@{}}
}{%
  \end{array}\right)
}

%--------grstep
% For denoting a Gauss' reduction step.
% Use as: \grstep{\rho_1+\rho_3} or \grstep[2\rho_5 \\ 3\rho_6]{\rho_1+\rho_3}
\newcommand{\grstep}[2][\relax]{%
   \ensuremath{\mathrel{
       {\mathop{\longrightarrow}\limits^{#2\mathstrut}_{
                                     \begin{subarray}{l} #1 \end{subarray}}}}}}
\newcommand{\swap}{\leftrightarrow}




\chead{\Large\textbf{Linear Transformations,  $\vec\S$1.8}}% 

\begin{document}
\thispagestyle{fancy}

Recall from a previous class:

\begin{itemize}

\item A set of vectors $\{ \vec{v}_1, \dots, \vec{v}_p\}$ in $\mathbb{R}^n$ is said to be {\bf linearly independent} if $$x_1 \vec{v}_1 +x_2 \vec{v}_2 +\dots +x_p \vec{v}_p = \vec{0}$$ has only the trivial solution. The set is called {\bf linearly dependent} if the equation has non-trivial solutions and the equation is called the {\bf dependence relation}.
\item Let $A= [\vec{a}_1 \; \dots \; \vec{a}_n ]$, then $\{ \vec{a}_1 ,  \dots , \vec{a}_n \}$ are linearly independent if and only if $A \vec{x} = \vec{0}$ has only the trivial solution.
\item {\bf Theorem 7:} A set $S=\{ \vec{v}_1, \dots, \vec{v}_p \}$ of two or more vectors is linearly dependent if and only if one of the vectors of $S$ is a linearly combination of the others.

Moreover if $\vec{v}_1\neq \vec{0}$, then some $\vec{v}_j$ for $j>1$ is a linear combination of $\vec{v}_1, \dots, \vec{v}_{j-1}$.
\item {\bf Theorem 8:} Any set $\{ \vec{v}_1, \dots, \vec{v}_p \}$ in $\mathbb{R}^n$ is linearly dependent if $p>n$.
\end{itemize}

\newpage

Define a transformation $T : \R^n \to \R^m$ to be a function from $\R^n$ to $\R^m$.   Recall we call $\R^n$ the domain and $\R^m$ the codmain. The set of outputs, $\{T(\vec{x}):\vec{x}\in \R^n\}$, is the image or the range of $T$.

One common example of a transformation $T$ is given by $T(\vec{x})=A\vec{x}$ where $A$ is an $m$ by $n$ matrix

\begin{defn}
A transformation $T$ is {\bf linear} if 
\begin{enumerate}
	\item $T(\vec{u}+\vec{v})=T(\vec{u})+T(\vec{v})$ for all $\vec{u},\vec{v}$ in $\R^n$,
	\item $T(c\vec{u})=cT(\vec{u})$ for all $\vec{u}$ in $\R^n$ and $c$ in $\R$.
\end{enumerate}
\end{defn}

\begin{enumerate}

\Problem {\bf Properties of Linear Transformations.}
\Part Prove if $T$ is linear, then $T(\vec{0})=\vec{0}$.

\vfill 

\Part $T(c\vec{u}+d\vec{v})=cT(\vec{u})+dT(\vec{v})$.

\vfill

\Problem Let $T(\vec{x})=A\vec{x}$ where $A=\begin{bmatrix}
1 & 0\\
-3 & 1\\
2 & -2\\
\end{bmatrix}$ and let $\vec{b}=\begin{bmatrix}
-2\\3\\-1\\
\end{bmatrix}$. 
\Part Is $\vec{b}$ in the range of $T$?

\vfill

\Part What is the range of $T$?

\vfill

\newpage

\Part Is $T$ surjective?

\vfill

\Part Is the transformation $T$ injective?

\vfill

\newpage


\Problem Let $T(\vec{x})=A\vec{x}$ where $A=\begin{bmatrix}
1 & 0 &-3\\
-3 & 1 & 6\\
2 & -2 & -1\\
\end{bmatrix}$ and let $\vec{b}=\begin{bmatrix}
-2\\3\\-1\\
\end{bmatrix}$. Is $\vec{b}$ in the range of $T$?

\vfill

\Problem Let $T:\R^n \to \R^m$ be linear and $\{\vec{v}_1, \vec{v}_2,\vec{v}_3\}$ be a linearly dependent set. Is $\{T(\vec{v}_1), T(\vec{v}_2),T(\vec{v}_3)\}$ linearly dependent?

\vfill


\Problem Let $T:\R^n \to \R^m$ be linear and $\{\vec{v}_1, \vec{v}_2,\vec{v}_3\}$ be a linearly \textbf{in}dependent set.  Is $\{T(\vec{v}_1), T(\vec{v}_2),T(\vec{v}_3)\}$ linearly \textbf{in}dependent?


%\Problem Suppose $T: \R^2 \to \R^3$ is linear and $

\vfill


\end{enumerate}

%\end{document}

\newpage
\chead{\Large\textbf{Linear Transformations -- Answers / Solutions}}%
\lhead{}
\rhead{}
\thispagestyle{fancy}

\begin{enumerate}
  \Answer 
  
  \Part Suppose $T$ is linear, then $$T(\vec{0}) = T(0 \cdot \vec{u})= 0 \cdot T(\vec{u})=\vec{0}.$$
  
  \Part Suppose $T$ is linear. By property (1) of linearity, we know $$T(c \vec{u}+d\vec{v}) =T(c\vec{u})+T(d\vec{v}).$$ By property (2), we know $$T(c\vec{u}))+T(d\vec{v})) =cT(\vec{u})+dT(\vec{v}).$$ Thus we have the result.
  
  \Answer
  
  \Part In order to determine if $\vec{b}$ is in the range of $T$, we solve the matrix equation $A\vec{x}=\vec{b}.$ Forming the augmented matrix and row-reducing, we get 
  $$\begin{amatrix}{2}
  1 & 0 & -2\\
  -3 & 1 & 3\\
  2 & -2 & -1\\
  \end{amatrix}\sim
  \begin{amatrix}{2}
  1 & 0 & -2\\
  0 & 1 & -3\\
  0 & 0 & -3\\
\end{amatrix}.$$

Thus there is no solution to $A\vec{x}=\vec{b}$.
  
  \Part Replacing $\vec{b}$ for the previous part with arbitrary components will give the range of $T$. We get 
  $$\begin{amatrix}{2}
  1 & 0 & b_1\\
  -3 & 1 & b_2\\
  2 & -2 & b_3\\
  \end{amatrix}\sim
  \begin{amatrix}{2}
  1 & 0 & b_1\\
  0 & 1 & b_2+3b_1\\
  0 & 0 & b_3+2b_2+4b_1\\
  \end{amatrix}.$$
  Thus $\vec{b}$ is in the range of $T$ if and only if $b_3+2b_2+4b_1=0$. We can see that the range is a plane in $\R^3$.
  
  \Part No, $T$ is not surjective. In particular, we showed in part (a) that there is no element in the domain that will map to $\vec{b}=\begin{bmatrix}
  -2\\3\\-1\\
  \end{bmatrix}$ under $T$. Equivalently we can say that $T(\vec{x})=\vec{b}$ has no solutions.
  
  \Part Recall $T$ is injective if for all $\vec{b}$, there is at most one $\vec{x}$ such that $T(\vec{x})=\vec{b}$. We use the contrapositive form of injectivity. Suppose $T(\vec{x})=T(\vec{y})$, then we get $A\vec{x}=A\vec{y}$. Rearranging the equation we get $A\vec{x}-A\vec{y}=\vec{0}$, and factoring, we get $A(\vec{x}-\vec{y})=\vec{0}$. We want to show that $A\vec{z}=\vec{0}$ has only the trivial solution. Forming the augmented matrix and row-reducing, we find $A\vec{z}=\vec{0}$ has only the trivial solution. Therefore $A(\vec{x}-\vec{y})=\vec{0}$ if and only if $\vec{x}-\vec{y}=\vec{0}$, and hence $\vec{x}=\vec{y}$. Thus $T$ is injective.
  
  \Answer Yes, $\vec{b}$ is in the range of $T$. In order to determine if $\vec{b}$ is in the range of $T$, we solve the matrix equation $A\vec{x}=\vec{b}.$ Forming the augmented matrix and row-reducing, we get 
  $$\begin{amatrix}{3}
  1 & 0 & -3 &-2\\
  -3 & 1 & 6 & 3\\
  2 & -2 & -1 & -1\\
  \end{amatrix}\sim
  \begin{amatrix}{3}
  1 & 0 & 0&7\\
  0 & 1 & 0&6\\
  0 & 0 & 1&3\\
  \end{amatrix}.$$
  
  Thus there is a unique solution to $A\vec{x}=\vec{b}$.
  
  \Answer Let $T:\R^n \to \R^m$ be linear and $\{\vec{v}_1, \vec{v}_2,\vec{v}_3\}$ be a linearly dependent set. Since $\{\vec{v}_1, \vec{v}_2,\vec{v}_3\}$ is linearly dependent, we know there exists $c_1, c_2, c_3$, not all zero, such that $c_1 \vec{v}_1+ c_2 \vec{v}_2+c_3\vec{v}_3 =\vec{0}.$ Applying $T$ to both sides of the equation we get $$T(c_1 \vec{v}_1+ c_2 \vec{v}_2+c_3\vec{v}_3) =T(\vec{0}).$$ Using linearity of $T$, we get $$c_1 T(\vec{v}_1)+ c_2T( \vec{v}_2)+c_3T(\vec{v}_3) =\vec{0}.$$ Since not all $c_1, c_2, c_3$ are zero, we have a non-trivial solution to the vector equation
  $$x_1 T(\vec{v}_1)+x_2 T(\vec{v}_2)+x_3 T(\vec{v}_3)=\vec{0},$$ and hence $\{T(\vec{v}_1), T(\vec{v}_2),T(\vec{v}_3)\}$ is linearly dependent.

\Answer No.  For example the transformation $T(\vec x) = A \vec x$ where $A$ is the $m$ by $n$ matrix with all entries equal to zero is a linear transformation.  However $\{T(\vec v_1), T(\vec v_2), T(\vec v_3)\}= \{ \vec 0, \vec 0, \vec 0\}$ is a linearly dependent set for any vectors $\vec v_i$. 


\end{enumerate}

\end{document}


%%% Local ables: 
%%% mode: latex
%%% TeX-master: t
%%% End: 
